% !TeX root = ../main.tex
% Draft from:
% - E:\GHuy\Study\TGMT\train_log.txt
% - E:\GHuy\Study\TGMT\transformer_best.pth
% - branch train (transreid_colab)
% - branch train: transreid_colab/README.md
% - branch train: transreid_colab/README_VG_LOCAL.md

\section{Thực nghiệm cho mô hình cải tiến trên nhánh train}
\label{sec:improved-train-experiment}

\subsection{Mô hình và hướng cải tiến}
Thực nghiệm này đánh giá một biến thể cải tiến của TransReID
\cite{he2021transreidtransformerbasedobjectreidentification}. So với cấu hình TransReID stride kết hợp JPM, mô hình
cải tiến được xây dựng theo bốn hướng chính:
\begin{itemize}
    \item \textbf{Visibility-guided patch weighting}: mỗi patch token được gán
    một trọng số nhìn thấy thông qua một đầu dự đoán hai lớp có kích thước
    $768 \rightarrow 192 \rightarrow 1$ kèm hàm sigmoid.
    \item \textbf{Dynamic top-k token enrichment}: các patch có trọng số cao
    nhất được dùng để tạo một \textit{visible prototype}, sau đó quay lại tăng
    cường cho toàn bộ patch token trước khi tổng hợp đặc trưng toàn cục.
    \item \textbf{JPM fusion}: đặc trưng cuối cùng ở pha suy luận được ghép từ
    nhánh toàn cục và bốn nhánh cục bộ của JPM nhằm giữ đồng thời thông tin
    toàn cục và chi tiết vùng.
    \item \textbf{Synthetic occlusion augmentation (OCC\_AUG)}: trong quá trình
    huấn luyện, hệ thống sinh thêm các mẫu bị che khuất tổng hợp để tăng khả
    năng chịu che khuất và ép đặc trưng của ảnh gốc và ảnh bị che vẫn giữ tính
    nhất quán.
\end{itemize}

Gọi $z_i$ là patch token thứ $i$, trọng số nhìn thấy tương ứng được ước lượng
theo công thức:
\begin{equation}
w_i = \sigma(\mathrm{MLP}(z_i)).
\end{equation}
Từ các patch có trọng số cao nhất, mô hình tạo ra một prototype:
\begin{equation}
p = \sum_{i \in \mathrm{TopK}} \alpha_i z_i, \qquad
\alpha_i = \frac{\exp(w_i)}{\sum_{j \in \mathrm{TopK}} \exp(w_j)}.
\end{equation}
Patch token sau tăng cường được tính như sau:
\begin{equation}
\tilde{z}_i = z_i + \lambda \, w_i \, p,
\end{equation}
trong đó $\lambda = 0.5$ theo cấu hình huấn luyện. Đặc trưng toàn cục dùng cho
truy hồi được tổng hợp bằng trung bình có trọng số:
\begin{equation}
g = \frac{\sum_i w_i \tilde{z}_i}{\sum_i w_i + \epsilon},
\end{equation}
với $\epsilon = 10^{-6}$.

Trong nhánh \texttt{TOKEN\_ENRICH}, tỉ lệ top-k được đặt là $0.3$. Với cấu hình
stride $[12,12]$ và ảnh đầu vào $256 \times 128$, lưới patch có kích thước
$21 \times 10$, tương ứng khoảng $210$ patch token; do đó nhánh token
enrichment sử dụng xấp xỉ $63$ patch có điểm nhìn thấy cao nhất để xây dựng
prototype. Với \texttt{OCC\_AUG}, một phần mẫu huấn luyện được tạo phiên bản bị
che khuất tổng hợp với xác suất $0.5$, diện tích vùng che trong khoảng
$[0.08, 0.30]$ và tỉ lệ khung che trong khoảng $[0.5, 2.0]$.

Về tối ưu hóa, mô hình vẫn kế thừa cơ chế \textit{softmax + soft triplet} của
TransReID trên nhánh toàn cục và các nhánh JPM. Trên đó, nhánh
\texttt{VIS\_WEIGHT} bổ sung một thành phần điều chuẩn cho patch weight:
\begin{equation}
\mathcal{L}_{reg}
= \alpha \, (\bar{w} - \tau)^2 - \beta \, \mathrm{Var}(w),
\end{equation}
trong đó $\tau = 0.5$, $\alpha = 1.0$, $\beta = 1.0$. Thành phần này khuyến
khích phân phối patch weight không bị suy biến về hằng số, đồng thời giữ giá trị
trung bình gần ngưỡng mục tiêu.

Khi bật \texttt{OCC\_AUG}, quá trình huấn luyện còn bổ sung loss nhận dạng trên
mẫu bị che khuất và một thành phần nhất quán đặc trưng giữa ảnh gốc và ảnh bị
che tổng hợp. Khi đó, loss tổng quát có thể được viết dưới dạng:
\begin{equation}
\mathcal{L}
= \mathcal{L}_{ID}^{JPM}
+ \mathcal{L}_{tri}^{JPM}
+ \lambda_{reg}\mathcal{L}_{reg}
+ \lambda_{occ}\mathcal{L}_{occ}
+ \lambda_{cons}\mathcal{L}_{cons},
\end{equation}
với $\lambda_{reg} = 0.05$, $\lambda_{occ} = 0.5$ và $\lambda_{cons} = 0.3$.
Trong đó, $\mathcal{L}_{occ}$ giúp mô hình vẫn giữ khả năng nhận dạng khi một
phần cơ thể bị che khuất, còn $\mathcal{L}_{cons}$ ép biểu diễn của cặp
``ảnh gốc -- ảnh bị che tổng hợp'' không lệch nhau quá xa trong không gian đặc
trưng.

\subsection{Thiết lập dữ liệu và cấu hình huấn luyện}
Thực nghiệm được thực hiện trên bộ dữ liệu Market-1501
\cite{zheng2015market1501}. Theo định nghĩa dataset trong nhánh
\texttt{train}, bộ dữ liệu được chia theo chuẩn gồm $12{,}936$ ảnh huấn luyện
cho $751$ định danh, $3{,}368$ ảnh query và $15{,}913$ ảnh gallery. Kích thước
ảnh đầu vào ở cả huấn luyện và đánh giá là $256 \times 128$.

Để làm rõ mức cải thiện, phần này sử dụng hai cấu hình tại cùng mốc
\textbf{epoch 120}:
\begin{itemize}
    \item \textbf{Baseline TransReID gốc}: giữ nguyên cấu hình TransReID stride
    trên Market-1501 với \texttt{vit\_base\_patch16\_224\_TransReID},
    \texttt{STRIDE\_SIZE=[12,12]}, \texttt{SIE\_CAMERA=True},
    \texttt{JPM=True}, đồng thời tắt toàn bộ các thành phần bổ sung gồm
    \texttt{VIS\_WEIGHT}, \texttt{TOKEN\_ENRICH}, \texttt{OCC\_AUG} và
    \texttt{LOCAL\_GROUP}.
    \item \textbf{Mô hình cải tiến}: giữ nguyên backbone TransReID stride và bật
    \texttt{VIS\_WEIGHT}, \texttt{TOKEN\_ENRICH} và \texttt{OCC\_AUG}, trong khi
    \texttt{LOCAL\_GROUP} được lược bỏ.
\end{itemize}

Ở pha suy luận, hệ thống sử dụng \texttt{USE\_JPM\_FUSION=True}, nghĩa là đặc
trưng truy hồi cuối cùng được ghép từ nhánh toàn cục và bốn nhánh cục bộ JPM.

Siêu tham số cho huấn luyện:
\begin{itemize}
    \item optimizer SGD, learning rate ban đầu $0.008$.
    \item batch size $64$, với \texttt{NUM\_INSTANCE=4}.
    \item weight decay $10^{-4}$.
    \item random horizontal flip với xác suất $0.5$.
    \item random erasing với xác suất $0.5$.
\end{itemize}

\subsection{Kết quả định lượng}

\begin{table}[t]
\centering
\caption{So sánh baseline TransReID gốc và mô hình cải tiến trên Market-1501.}
\label{tab:baseline-vs-improved-epoch120}
\begin{tabular}{l|c|c|c}
\hline
Mô hình & Epoch & mAP (\%) & Rank-1 (\%) \\
\hline
Baseline & 120 & 88.9 & 95.2 \\
Mô hình cải tiến & 120 & 89.3 & 95.1 \\
\hline
\end{tabular}
\end{table}

So với baseline TransReID gốc, mô hình cải tiến tăng khoảng $0.3$ điểm mAP,
trong khi Rank-1 gần như giữ nguyên. Xu hướng này phù hợp với mục tiêu của
\textit{visibility-guided patch weighting}, \textit{token enrichment} và
\textit{OCC\_AUG}: cải thiện chất lượng biểu diễn và thứ tự truy hồi tổng thể,
nhiều hơn là thay đổi mạnh top-1 accuracy.
